\documentclass[10pt,a4paper]{amsart}
\usepackage[margin=19mm]{geometry}
\usepackage{amsmath,amssymb,mathtools}
\usepackage[scr=rsfs,cal=euler]{mathalfa}
\usepackage{xfrac}
\usepackage[unicode,hidelinks,bookmarks=false]{hyperref}
\newcommand{\ie}{i.e.\ }
\newcommand{\cf}{cf.\ }
\newcommand{\resp}{respectively\ }
\newcommand{\Subsection}{\subsection}
\providecommand{\nobreakdash}{\nobreak\hbox{-}}
\setlength{\emergencystretch}{2em}
\multlinegap0pt
\usepackage{bm}
\usepackage{bbm}
\usepackage{dsfont}
\usepackage{xspace}
\usepackage{cancel}
\usepackage{mathtools}
\usepackage{amsthm}
\makeatletter
\@ifundefined{X}{\newtheorem{X}{X}}{}
\@ifundefined{lemma}{\newtheorem{lemma}[X]{Lemma}}{}
\makeatother
\newcommand{\ZZ}{\ensuremath{\mathbb{Z}}}
\newcommand{\cL}{\ensuremath{\mathcal{L}}}
\renewcommand{\le}{\ensuremath{\leqslant}}
\title[Exponential sums with polynomials and their applications to ]{Exponential sums with polynomials and their applications to primes in sparse sets}
\author{Lingyu Guo, Victor Zhenyu Guo, Mengyao Jing}
\date{Source: arXiv:2510.20562v1 (CC BY 4.0)}
\begin{document}
\maketitle

\noindent\emph{Source and licence.} Exponential sums with polynomials and their applications to primes in sparse sets. Version arXiv:2510.20562v1 (CC BY 4.0), distributed under the Creative Commons Attribution 4.0 International licence, as recorded in the arXiv licence metadata for this exact version and shown on the abstract page https://arxiv.org/abs/2510.20562v1. Full rights notice: licenses/arxiv-2510.20562-CC-BY-4.0.txt.\par\smallskip

\noindent\textit{Editorial scope.} Counted unit: the Lemma whose source label is \texttt{lemma: Gamma II} in the article (source0.tex lines 1743--1753, source label \texttt{lemma: Gamma II}), delivered together with the article's own complete original proof of it (source0.tex lines 1755--1877, one \texttt{proof} environment of 123 lines). No mathematics is written, repaired or re-derived: statement and proof are the article's own text line for line. Required context retained uncounted: lemma: zhai's lemma (lines 562--573); lem: exponent pair method (lines 577--588); lemmea: balance (lines 640--651); lem:A (lines 659--665). Every definition, display and label of those lines is retained unchanged. Omitted from this package and not redistributed: 1--561, 574--576, 589--639, 652--658, 666--1742, 1754--1754, 1878--2142. Nothing inside a retained excerpt is omitted or altered: every retained body is delivered exactly as the article's lines are written, so no omission receipt of any kind accompanies this package. Citations: the retained excerpts carry no citation command at all, so no bibliography entry of the article is transcribed and no reference is redistributed. Reference adaptation: none was needed and none was made; every reference inside the delivered unit resolves inside it, so no unresolved reference or citation is produced. Printed slips kept literal: no printed word, symbol, hypothesis or number of the counted statement or of its proof was altered. Layout and class: the article's own class is \texttt{amsart} with packages including amsmath, amssymb, amsbsy, amsfonts, amsthm, latexsym, amsopn, amstext, amsxtra, euscript, amscd, stmaryrd, mathrsfs, cite; the wrapper is the lane's pinned \texttt{amsart} editorial wrapper at 10pt on A4, which restates the theorem-like environments the retained text needs and adds the article's own macro definitions that the retained text needs (\texttt{\textbackslash ZZ}, \texttt{\textbackslash cL}, \texttt{\textbackslash le}), with extra wrapper packages (bm, bbm, dsfont, xspace, cancel, mathtools). The delivered numbering is the unit's own. Assets: no retained span contains a figure environment, a graphics-inclusion command, a PDF-page inclusion, a table or a TikZ picture; no image file is copied into this package, the package declares no asset dependency and no image file is redistributed with it. Licence: version arXiv:2510.20562v1 of this article is distributed under the Creative Commons Attribution 4.0 International licence, as recorded in the arXiv licence metadata for this exact version and shown on the abstract page https://arxiv.org/abs/2510.20562v1. A full rights notice is delivered at licenses/arxiv-2510.20562-CC-BY-4.0.txt. Characters: every retained excerpt is pure ASCII, so the delivered engine, XeLaTeX with its default Latin Modern text font, reports no missing character.
\begin{lemma}\label{lemma: zhai's lemma}
Let $d\geqslant 2$ be a fixed integer. Suppose that $a_1,\cdots,a_d$ are non-zero real numbers,
$\alpha_1,\cdots,\alpha_d$ are distinct real constants such that $\alpha_1,\cdots,\alpha_d\not\in\mathbb{Z}$,
and assume that $\Theta>0$. Let $\mathcal{I}\subset[1,2]$ be a subinterval such that
$$
\left|a_1 t^{\alpha_1} + \cdots+a_d t^{\alpha_d} \right|\leqslant \Theta,\ \ \ t\in\mathcal{I}.
$$
Then we have
$$
|\mathcal{I}|\ll \left(\frac{\Theta}{|a_1|+\cdots+|a_d|}\right)^{1/(d-1)}.
$$
\end{lemma}

\begin{lemma}\label{lem: exponent pair method}
Let $s\geqslant 2$ be a positive integer, and let $f(x)$ be a real valued
function with $s$ continuous derivatives on $[N,2N]$ such that
$$
|f^{(s)}(n)|\asymp YN^{1-s}.
$$
Then we have
$$
\sum_{n\sim N}\mathbf{e}(f(n))\ll Y^{\kappa}N^{\lambda}+Y^{-1},
$$
where $(\kappa,\lambda)$ is any exponent pair.
\end{lemma}

\begin{lemma} \label{lemmea: balance} Let
$$
L(E)=\sum_{1\leqslant i\leqslant u}A_iE^{a_i}+\sum_{1\leqslant j\leqslant v}B_jE^{-b_j},
$$
where $A_i,B_j,a_i$ and $b_j$ are positive. Let $0\leqslant E_1\leqslant E_2$. Then there is
some $E\in(E_1,E_2]$ such that
$$
L(E)\ll \sum_{1\leqslant i\leqslant u}\sum_{1\leqslant j\leqslant v}\left(A_i^{b_j}B_j^{a_i}\right)^{1/(a_i+b_j)}
+\sum_{1\leqslant i\leqslant u}A_iE_1^{a_i}+\sum_{1\leqslant j\leqslant v}B_jE_2^{-b_j},
$$
where the implied constant only depends on $u$ and $v$.
\end{lemma}

\begin{lemma}
\label{lem:A}
Suppose $f(n)$ is a complex valued function and $I$ is an interval such that $f(n) = 0$ if $n \notin I$. If $Q$ is a positive integer then
$$
|\sum_{n \in I} f(n)|^2 \le \frac{|I| + Q}{Q} \sum_{|q|< Q} \big(1 - \frac{q}{Q} \big) \sum_{n \in I} f(n) \overline{f(n-q)}.
$$
\end{lemma}

\begin{lemma}\label{lemma: Gamma II} Assume that $X^{1/2}<M<X^{2/3}$, $W=X^{1-\gamma_1\gamma_2+\varepsilon}$, $j<W^{1+2\varepsilon}$ and
	\begin{align}\label{eq: condition 3}
		\left\{\begin{array}{ll} 26\gamma_1\gamma_2-2\gamma_2-23>0, & \\
			15\gamma_1\gamma_2-14>0,
		\end{array}\right.
	\end{align}
	we have
	\begin{align*}
		\sum_{\substack{(k,h) \in \ZZ^2 \backslash (0,0) \\ |k|<\cL X^{1-\gamma_2} \\ |h|<\cL X^{\gamma_2(1-\gamma_1)}}}|\mathfrak{T}_{II}|\ll X^{\gamma_1\gamma_2-2\varepsilon}.
	\end{align*}
\end{lemma}

\begin{proof}
	We divide the range $m\sim M$ into the following two parts:
	$$
	\mathcal{J}_3=\bigg\{m\sim M:
	\bigg|hX^{\gamma_1\gamma_2}
	\Big(\frac{m}{M}\Big)^{\gamma_1\gamma_2}
	+(k+j)X^{\gamma_2}\Big(\frac{m}{M}\Big)^{\gamma_2}\bigg|
	<X^{12+\gamma_2-13\gamma_1\gamma_2+24\varepsilon}\bigg\}
	$$
	and
	$$
	\mathcal{J}_4=\{m\sim M: m\notin\mathcal{J}_3\}.
	$$
	Then we have
	\begin{align*}
		\mathfrak{T}_{II}=\mathfrak{T}_{II,1}+\mathfrak{T}_{II,2},
	\end{align*}
	where
	$$
	\mathfrak{T}_{II,1}=\mathop{\sum_{m\sim M}}_{m\in\mathcal{J}_3}\mathop{\sum_{n\sim N}}_{X < mn \leqslant aX}
	a_mb_n\mathbf{e}\Big(h\Big((mn+1)^{\gamma_2}-\frac{w}{W}\Big)^{\gamma_1}
	+k(mn)^{\gamma_2}+j(mn+1)^{\gamma_2}\Big),
	$$
	and
	$$
	\mathfrak{T}_{II,2}=\mathop{\sum_{m\sim M}}_{m\in\mathcal{J}_4}\mathop{\sum_{n\sim N}}_{X < mn \leqslant aX}
	a_mb_n\mathbf{e}\left(h\left((mn+1)^{\gamma_2}-\frac{w}{W}\right)^{\gamma_1}
	+k(mn)^{\gamma_2}+j(mn+1)^{\gamma_2}\right).
	$$
	
	We begin with $\mathfrak{T}_{II,1}$. Clearly,
	\begin{align*}
		\mathfrak{T}_{II,1}\ll X^{2\varepsilon}N\mathop{\sum_{m\sim M}}_{m\in\mathcal{J}_3}1
		\ll X^{2\varepsilon}N\#\mathcal{J}_3.
	\end{align*}
	It follows from Lemma \ref{lemma: zhai's lemma} that
	\begin{align*}
		\#\mathcal{J}_3\ll \bigg(\frac{X^{12+\gamma_2-13\gamma_1\gamma_2+24\varepsilon}}
		{|hX^{\gamma_1\gamma_2}|
			+|(k+j)X^{\gamma_2}|}\bigg)M.
	\end{align*}
This can be bounded by considering two cases: $h\neq0$ with $j+k=0$; or $h=0$
   with $j+k\neq0$. By $15\gamma_1\gamma_2-14>0$ we have
	\begin{align*}
		\sum_{\substack{(k,h) \in \ZZ^2 \backslash (0,0) \\ |k|<\cL X^{1-\gamma_2} \\ |h|<\cL X^{\gamma_2(1-\gamma_1)}}}|\mathfrak{T}_{II,1}|
		\ll X^{\gamma_1\gamma_2-2\varepsilon}.
	\end{align*}
	
	Now we estimate $\mathfrak{T}_{II,2}$.
	Applying the Cauchy inequality we have
	\begin{align*}
		|\mathfrak{T}_{II,2}|^2\leqslant X^{\varepsilon}M\sum_{m\sim M}
		\Bigg|\sum_{n\sim N}b_n\mathbf{e}&\bigg(\left(h(mn+1)^{\gamma_2}-\frac{w}{W}\right)^{\gamma_1}\\
		&+k(mn)^{\gamma_2}+j(mn+1)^{\gamma_2}\bigg)\Bigg|^2.
	\end{align*}
	Let $Q\leqslant N$ be an integer, to be chosen later. The Weyl-van der Corput inequality
   (see Lemma \ref{lem:A}) yields
	\begin{align*}
		|\mathfrak{T}_{II,2}|^2\ll\frac{X^{2+2\varepsilon}}{Q}
		+\frac{X^{1+\varepsilon}}{Q}\sum_{1\leqslant |q|\leqslant Q}
		\sum_{n\sim N}\Bigg|\sum_{m\sim M}\mathbf{e}(f(m))\Bigg|,
	\end{align*}
	where
	\begin{align*}
		f(x)=&h\left((xn+xq+1)^{\gamma_2}-\frac{w}{W}\right)^{\gamma_1}
		+k(xn+xq)^{\gamma_2}+j(xn+xq+1)^{\gamma_2}\\
		&-h\left((xn+1)^{\gamma_2}-\frac{w}{W}\right)^{\gamma_1}
		+k(xn)^{\gamma_2}+j(xn+1)^{\gamma_2}.
	\end{align*}
	For fixed $n\sim N$ and $q$, and for $x\sim M$, it is easy to see that
	\begin{align*}
		f^{(s)}(x)\asymp\Big|&(\gamma_1\gamma_2)(\gamma_1\gamma_2-1)\cdots(\gamma_1\gamma_2-s+1)
		hX^{\gamma_1\gamma_2}\\
		&+\gamma_2(\gamma_2-1)\cdots(\gamma_2-s+1)(j+k)X^{\gamma_2}\Big||q|X^{-1}M^{1-s}.
	\end{align*}
	
	Let
	$$
	\Delta_4=\left|hX^{\gamma_1\gamma_2}
	\left(\frac{m}{M}\right)^{\gamma_1\gamma_2}
	+(k+j)X^{\gamma_2}\left(\frac{m}{M}\right)^{\gamma_2}\right|.
	$$
	The condition $m\in\mathcal{J}_4$ implies that
	$$
   X^{12+\gamma_2-13\gamma_1\gamma_2+24\varepsilon}\ll \Delta_4\ll
   X^{1+\gamma_2-\gamma_1\gamma_2+2\varepsilon}.
   $$
	Applying Lemma \ref{lem: exponent pair method} with $(\kappa,\lambda)=(1/2,1/2)$
	we can see that
	\begin{align*}
		\mathfrak{T}_{II,2}\ll X^{(3-\gamma_1\gamma_2+\gamma_2)/4+\varepsilon}N^{1/4}Q^{1/4}
		+X^{(13\gamma_1\gamma_2-\gamma_2-10)/2-11\varepsilon}N^{1/2}Q^{-1/2}
		+X^{1+\varepsilon}Q^{-1/2}.
	\end{align*}
	This bound is trivial when $Q<1$, so we may restrict that $0<Q\leqslant N$.
	Lemma \ref{lemmea: balance} and $X^{1/3}<N<X^{1/2}$ yield
    \begin{align*}
    \mathfrak{T}_{II,2}\ll
		X^{2\gamma_1\gamma_2-1-3\varepsilon}
       +X^{(11-2\gamma_1\gamma_2+2\gamma_2)/12+\varepsilon}
       +X^{(13\gamma_1\gamma_2-\gamma_2-10)/2-11\varepsilon}
       +X^{5/6+\varepsilon}.
    \end{align*}
    It follows from the inequalities that
    $$
    26\gamma_1\gamma_2-2\gamma_2-23>0\ \ \ \hbox{and}
    \ \ \ 12\gamma_1\gamma_2-11>0
    $$
    we have
	\begin{align*}
		\sum_{\substack{(k,h) \in \ZZ^2 \backslash (0,0) \\ |k|<\cL X^{1-\gamma_2} \\ |h|<\cL X^{\gamma_2(1-\gamma_1)}}}|\mathfrak{T}_{II,2}|
     \ll X^{\gamma_1\gamma_2-2\varepsilon}.
	\end{align*}
Here we also remark that $12\gamma_1\gamma_2>11$ is derived from $15\gamma_1\gamma_2>14$.

   Therefore,
   $$
   \sum_{\substack{(k,h) \in \ZZ^2 \backslash (0,0) \\ |k|<\cL X^{1-\gamma_2} \\ |h|<\cL X^{\gamma_2(1-\gamma_1)}}}|\mathfrak{T}_{II}|
   \ll \sum_{\substack{(k,h) \in \ZZ^2 \backslash (0,0) \\ |k|<\cL X^{1-\gamma_2} \\ |h|<\cL X^{\gamma_2(1-\gamma_1)}}}|\mathfrak{T}_{II,1}|
   +\sum_{\substack{(k,h) \in \ZZ^2 \backslash (0,0) \\ |k|<\cL X^{1-\gamma_2} \\ |h|<\cL X^{\gamma_2(1-\gamma_1)}}}|\mathfrak{T}_{II,2}|
     \ll X^{\gamma_1\gamma_2-2\varepsilon}.
   $$
\end{proof}

\end{document}
